\chapter{Zipf's Law and Token Distributions}
\label{ch16}

Empirical linguistics has known for nearly a century that word frequencies in large corpora are not spread evenly across the vocabulary. A handful of words dominate, while the vast majority appear rarely. This chapter makes that observation precise, reviews the main theoretical explanations (Pitman--Yor processes, combinatorial growth arguments, and the quantization model), and proposes how the bivariate generating function machinery developed in Chapters~4--5 might be brought to bear on the problem.

\section{Zipf's Law}
\label{ch16:zipf-law}

\begin{definition}[Zipf's Law]
\label{ch16:zipf-def}
Let a corpus of tokens be given. Rank the distinct token types by decreasing frequency, so type ranked $r$ is the $r$-th most common. Zipf's law asserts that the frequency of the $r$-th type is approximately
\[
f(r) \;\approx\; \frac{C}{r^{s}},
\]
where $C > 0$ is a corpus-dependent constant and the exponent $s \approx 1$ for natural language corpora.
\end{definition}

This is an \emph{empirical observation}, not a theorem. The exponent $s$ is close to 1 across many languages and many corpora, which is striking, but it is not a mathematical law with a proof. Measuring $s$ requires fitting a power law to noisy rank-frequency data, which introduces its own methodological complications (see Piantadosi's review below).

The connection to the complementary cumulative distribution is direct: if $f(r) \propto r^{-s}$, then the probability that a randomly drawn token has rank at most $r$ grows as $\sum_{k=1}^{r} k^{-s}$, and by Euler--Maclaurin the probability that a token's frequency exceeds some threshold $f$ satisfies $\Pr[\text{freq} \ge f] \propto f^{-1/s}$. For $s=1$ this gives a power-law tail with exponent $1$, sometimes written as saying the frequency distribution has a Pareto tail of index $1$.

\subsection{Mandelbrot's Refinement}
\label{ch16:mandelbrot}

Mandelbrot observed that the simple formula $C/r^{s}$ fits poorly for small ranks (the very top of the frequency list) and proposed
\[
f(r) \;\approx\; \frac{C}{(r + b)^{s}},
\]
where $b \ge 0$ is a shift parameter. This one-parameter extension gives noticeably better empirical fits, particularly for high-frequency function words. Mandelbrot connected this formula to information-theoretic arguments about optimal coding, though those arguments remain contested. For our purposes, the Mandelbrot form is a useful phenomenological model; it does not by itself explain \emph{why} such a distribution arises.

\section{Piantadosi's Critical Review}
\label{ch16:piantadosi}

The apparent universality of Zipf's law has attracted an enormous variety of proposed explanations: preferential attachment (Simon 1955), optimal coding (Mandelbrot), maximum-entropy arguments, and more. Piantadosi \cite{Piantadosi2014} provides a systematic critical evaluation of these explanations.

The central point of \cite{Piantadosi2014} is methodological: virtually every proposed generative model for language---from random text generated by a monkey typing uniformly on a keyboard, to preferential-attachment processes, to context-free grammars---can be tuned to produce a power-law rank--frequency distribution. Matching Zipf's law is therefore \emph{weak evidence} for any particular explanation. A model that predicts Zipf's law is not confirmed; it has merely failed to be ruled out by this one summary statistic.

Piantadosi argues that a more informative test would involve the \emph{joint} distribution of token frequency with other observables---word length, syntactic category, contextual predictability---rather than frequency alone. For example, there is a well-documented negative correlation between word frequency and word length (high-frequency words tend to be short), and any serious explanation of Zipf's law should also account for this. Models that reproduce the rank--frequency plot while getting the length distribution wrong are incomplete.

This critical stance is useful for us: it tells us that finding a generating function model whose singularity analysis yields a $1/r$ rank distribution is not, by itself, evidence that the model is correct. The bivariate generating function proposal in Section~\ref{ch16:bgf} is motivated partly by the desire to track \emph{both} rank and length simultaneously, as Piantadosi recommends.

\section{Pitman--Yor Processes and Power Laws}
\label{ch16:py}

The most mathematically clean explanation of Zipf's law in the language modeling literature is via the Pitman--Yor (PY) process \cite{GoldwaterGriffithsJohnson2011}.

\subsection{The Pitman--Yor Process}
\label{ch16:py-def}

The Pitman--Yor process with discount parameter $a \in (0,1)$ and concentration parameter $\theta > -a$ is a distribution over infinite sequences of tokens. It can be described by the \emph{Chinese restaurant process}: tokens arrive one at a time. When the $n$-th token arrives and $k$ distinct types have appeared so far, with type $i$ having count $c_i$:
\begin{itemize}
\item with probability $(c_i - a)/(n - 1 + \theta)$, emit type $i$ (reinforcement),
\item with probability $(\theta + ak)/(n - 1 + \theta)$, emit a brand-new type (innovation).
\end{itemize}

The key property is that the marginal frequency distribution of types, in the limit of a long sequence, has a power-law tail. Specifically \cite{GoldwaterGriffithsJohnson2011}:
\[
\Pr[\text{$r$-th most frequent type after $n$ tokens}] \;\sim\; C \, n^{-a} \cdot r^{-(1+1/a)},
\]
so the rank--frequency exponent is $s = 1 + 1/a$. For $a = 0.5$, we get $s = 3$; for $a$ close to $1$, $s$ approaches $2$ from above; for $a$ close to $0$, $s$ grows without bound (the distribution becomes more concentrated). To recover empirical Zipf with $s \approx 1$, one would need $a$ close to $1$.

\begin{remark}
\label{ch16:py-remark}
The exponent $s = 1+1/a$ is greater than $1$ for all $a \in (0,1)$. So a single-level Pitman--Yor process cannot directly produce the empirical $s \approx 1$ unless $a \to 1$, which is a degenerate limit. The two-stage model below addresses this.
\end{remark}

\subsection{The Two-Stage Model}
\label{ch16:two-stage}

Goldwater, Griffiths, and Johnson \cite{GoldwaterGriffithsJohnson2011} propose a hierarchical (two-stage) construction: a \emph{base} Pitman--Yor process generates an abstract ``cached vocabulary,'' and a second-level process selects from this vocabulary or creates novel items. This composition of two PY processes allows the effective exponent to be lower and more flexibly tuned to match empirical data. The two-stage structure is natural from a linguistic standpoint: speakers draw on a mental lexicon (the cache), and the lexicon itself grows over time via a process with power-law dynamics.

\section{Berman's Combinatorial Derivation}
\label{ch16:berman}

A different perspective comes from purely combinatorial reasoning. Berman \cite{Berman2025} (arXiv:2511.17575) derives Zipf-type exponents without Bayesian priors or preferential attachment, working instead from the growth rate of the number of distinct types as a function of corpus size.

The argument proceeds roughly as follows. Consider a random text generated by a memoryless source. Track the \emph{type count} $V(n)$, the number of distinct token types seen after $n$ tokens. If $V(n) \sim n^{\beta}$ for some $\beta \in (0,1)$, a combinatorial counting argument on the expected multiplicity of a type at rank $r$ yields a Zipf exponent that is determined by $\beta$. The key insight is that the Zipf exponent is not a free parameter to be fit; it is constrained by the vocabulary growth exponent, which is in principle observable independently. This gives a combinatorial prediction: if you know how fast new words appear in a corpus, you can predict the shape of the rank--frequency plot.

This is qualitatively different from the PY explanation: Berman's derivation is deterministic in the sense that the exponent emerges from a combinatorial growth law, not from the random branching structure of a stochastic process. Whether the two explanations are ultimately equivalent (perhaps the PY vocabulary growth has $V(n) \sim n^{a}$, matching the PY discount parameter) is an open question worth investigating.

\section{Mikhaylovskiy's Temperature Window}
\label{ch16:mikhaylovskiy}

Chapter~15 established that autoregressive language models with temperature $T$ sampling undergo phase transitions in their output statistics as $T$ varies. Mikhaylovskiy \cite{Mikhaylovskiy2025b} connects this directly to Zipf's law.

The finding is: for transformer language models, there exists a finite interval $[T_{\min}, T_{\max}]$ of temperatures within which the empirical rank--frequency exponent of model outputs is approximately $1$ (matching natural language). Outside this interval, the exponent drifts---below $T_{\min}$, the model becomes repetitive and the distribution concentrates on a few high-frequency types, steepening the rank--frequency curve; above $T_{\max}$, the model becomes diffuse and the distribution flattens.

Mikhaylovskiy frames this as a \emph{universality window}: the Zipf behavior is not a property of any single temperature but of a region in parameter space, and the boundaries of this region are the two phase transitions identified in Chapter~15. This is conceptually important. It means that a transformer model does not ``naturally'' produce Zipf distributions; it requires that the temperature parameter be within a specific range. A poorly calibrated model (or a model sampled at $T=1$ without any tuning) may or may not fall inside this window.

\begin{remark}
\label{ch16:mikhaylovskiy-remark}
This also has a practical implication. If a practitioner observes that a language model's output does not follow Zipf's law, the diagnosis may not be a model failure but a temperature miscalibration. Conversely, observing Zipf-like output is not evidence that the model is generating ``natural'' text; it only says the temperature is inside the window.
\end{remark}

\section{The Quantization Model}
\label{ch16:quantization}

Michaud et al. \cite{MichaudEtAl2023} (arXiv:2303.13506) propose that the internal representations of a neural language model are organized around a discrete set of ``quanta''---atomic skills or features, each associated with a subset of the vocabulary. The model's outputs are produced by composing these quanta, and the usage frequency of each quantum follows a power law.

The quantization model makes two related claims:
\begin{enumerate}
\item The power-law shape of neural scaling laws (loss decreases as $L \sim N^{-\alpha}$ in parameters $N$) follows from the discrete-quanta structure: as model capacity increases, more quanta are acquired, each contributing a roughly equal increment of capability.
\item The Zipf-like frequency distribution of output tokens follows from the same structure: tokens associated with common quanta are produced frequently, tokens associated with rare quanta are produced rarely, and if quantum usage frequencies are themselves power-law distributed, token frequencies inherit this distribution.
\end{enumerate}

The connection to the Goldwater--Griffiths--Johnson two-stage model \cite{GoldwaterGriffithsJohnson2011} is direct: the quanta of Michaud et al. play the role of the ``cached vocabulary'' in the Pitman--Yor two-stage construction. A token is produced by first selecting a quantum (the first stage, PY-like) and then emitting a token associated with that quantum (the second stage). If quantum selection follows a power law with exponent determined by $a$, and each quantum maps to a roughly fixed number of tokens, the output token distribution inherits the power-law shape. The quantization model thus gives a mechanistic neural network interpretation of the abstract PY hierarchy.

\section{A Formal Proposition on Multinomial Models}
\label{ch16:proposition}

Before turning to the bivariate generating function proposal, we record a basic proposition that clarifies what Zipf's law does and does not require.

\begin{proposition}
\label{ch16:multinomial-prop}
Let tokens be drawn i.i.d. from a fixed vocabulary $\{w_1, w_2, \ldots\}$ with probabilities $p_1 \ge p_2 \ge \cdots > 0$. After $n$ draws, let $\hat{f}(r,n)$ denote the empirical frequency of the type with the $r$-th largest count. Then, as $n \to \infty$, $\hat{f}(r,n) \to p_r$ almost surely. The empirical rank--frequency plot converges to the true parameter curve $r \mapsto p_r$. Consequently, the empirical distribution is Zipfian with exponent $s$ if and only if $p_r \propto r^{-s}$.
\end{proposition}

\begin{proof}
By the strong law of large numbers, the empirical frequency of each type $w_i$ converges almost surely to $p_i$. Since the ranking is determined by these frequencies, and the $p_r$ are distinct (generically), the rank of each type stabilizes almost surely after finitely many steps. Therefore $\hat{f}(r,n) \to p_r$ almost surely as $n \to \infty$.
\end{proof}

\begin{remark}
\label{ch16:multinomial-remark}
This proposition says that a multinomial model with Zipfian parameters produces Zipfian output, and a multinomial model with non-Zipfian parameters does not. The proposition is trivial in one direction---it simply says empirical frequencies converge---but its point is conceptual: Zipf's law in the output is equivalent to Zipf's law in the model parameters when the model is i.i.d. The interesting scientific question is not whether an i.i.d. model can produce Zipf's law (it can, trivially, if parameterized correctly), but where the Zipfian parameter values come from. The PY, combinatorial, and quantization frameworks are answers to this deeper question.
\end{remark}

\section{Worked Mini-Example}
\label{ch16:example}

\begin{example}[Pitman--Yor with $a = 0.5$]
\label{ch16:py-example}
Set $a = 0.5$. The Pitman--Yor prediction is $s = 1 + 1/a = 1 + 2 = 3$. The rank--frequency plot should follow $f(r) \propto r^{-3}$, which is \emph{steeper} than empirical Zipf ($s \approx 1$). In a log--log plot, the slope would be $-3$ rather than $-1$. To achieve $s \approx 1$, one would need $a \approx 1$, but the PY process is only defined for $a < 1$. The two-stage construction of \cite{GoldwaterGriffithsJohnson2011} provides a way around this by composing two levels.
\end{example}

\begin{example}[Memoryless unigram model]
\label{ch16:unigram-example}
Consider a memoryless (i.i.d.) model with vocabulary $\{w_1, \ldots, w_V\}$ and uniform probabilities $p_i = 1/V$. By Proposition~\ref{ch16:multinomial-prop}, the empirical rank--frequency plot converges to the flat line $f(r) = 1/V$. This is \emph{not} Zipfian; it has exponent $s = 0$. If instead $p_i \propto i^{-1}$ (a manually specified Zipfian distribution), then the model outputs Zipfian frequencies, but we have simply assumed what we wanted to explain. The conclusion is that Zipf's law cannot emerge from a simple i.i.d. model unless it is pre-loaded into the parameters. This is why the process-level explanations---PY, combinatorial growth, quantization---are necessary: they explain how Zipfian parameters arise from simpler primitives.
\end{example}

These two examples illustrate a general point. Zipf's law is a \emph{non-trivial} statistical regularity that discriminates between model classes. An i.i.d. model with a fixed vocabulary will not generate it spontaneously. A model with memory, preferential attachment, or hierarchical structure can generate it, and the exponent encodes something about the structure of that memory or hierarchy.

\section{A Bivariate Generating Function Proposal}
\label{ch16:bgf}

The tools developed in Chapters~4--5 and formalized in Flajolet and Sedgewick \cite{FlajoletSedgewick2009}, Chapter~IX, handle bivariate generating functions (BGFs) of the form $F(z,u) = \sum_{n,k} f_{n,k} z^n u^k$, where the variable $u$ marks a secondary parameter. The coefficient $[z^n u^k] F(z,u)$ counts objects of size $n$ with secondary parameter value $k$.

We propose the following BGF as a tool for studying the joint distribution of token length and rank:
\[
F(z, u) \;=\; \sum_{w \in \mathcal{V}} z^{|w|} \, u^{\text{rank}(w)},
\]
where $\mathcal{V}$ is the vocabulary, $|w|$ is the character-length of token $w$, and $\text{rank}(w)$ is the frequency rank of $w$ in the corpus. If the rank--length joint distribution follows Zipf's law with a length--rank correlation, then $F(z,u)$ encodes this correlation in its singular behavior.

Specifically, the singularity of $F(z,u)$ in $z$, as a function of $u$, would shift as $u$ varies. If the dominant singularity $\rho(u)$ satisfies $\rho(u) \sim \rho_0 \cdot u^{\gamma}$ near $u = 1$, then by the transfer lemma of singularity analysis (Flajolet--Sedgewick Ch.~VI), the coefficient $[z^n] F(z,u)$ is dominated by tokens near a characteristic rank $r^\star(n) \sim n^{1/\gamma}$. Inverting, this would give $s = 1/\gamma$.

This is a proposal, not a theorem. Several technical difficulties stand in the way:
\begin{enumerate}
\item The vocabulary is finite in practice, so $F(z,u)$ is a polynomial in $z$ for each fixed $u$; a singularity analysis requires a suitable infinite or asymptotic model.
\item The rank function depends on the corpus, which is itself random; incorporating this randomness into the generating function requires a random analytic function framework.
\item The connection between the PY process or Berman's growth model and a specific BGF functional equation is not yet established.
\end{enumerate}

This is related to the open problems of Chapter~18. The nearest solved analogues are the BGF analyses of random tries (Flajolet--Sedgewick Ch.~IX), where length and depth play roles analogous to length and rank here. A systematic translation of the trie analysis to the token-rank setting would be a concrete step toward a theorem.

\section{Notes and Further Reading}
\label{ch16:notes}

Zipf's original data and the rank--frequency observation appear in \cite{Zipf1949} (not cited as a primary source here but standard background). Mandelbrot's information-theoretic derivation is in his 1953 paper. The Simon (1955) preferential-attachment model is the discrete-time precursor to Barabási--Albert networks and reproduces Zipf's law from a simple urn process.

Piantadosi \cite{Piantadosi2014} is the clearest statement of why reproducing the rank--frequency plot is insufficient evidence for a mechanistic model, and should be read alongside any of the theoretical derivations. Goldwater, Griffiths, and Johnson \cite{GoldwaterGriffithsJohnson2011} give the full Pitman--Yor derivation with empirical validation on English and Chinese corpora.

Berman \cite{Berman2025} (arXiv:2511.17575) provides the combinatorial perspective. Mikhaylovskiy \cite{Mikhaylovskiy2025b} provides the empirical temperature-window analysis for transformers. The quantization model of Michaud et al. \cite{MichaudEtAl2023} (arXiv:2303.13506) connects Zipf's law to neural scaling laws via a discrete-quanta structure.

For the BGF methodology, Flajolet and Sedgewick \cite{FlajoletSedgewick2009} Chapter~IX is the standard reference. Random tries are the worked-out example closest to the token-rank problem.

\subsection*{Recommended Reading}

\begin{itemize}
  \item \textbf{G.~K. Zipf}, \emph{Human Behavior and the Principle of Least Effort} (Addison-Wesley, 1949) \cite{Zipf1949} --- The original empirical study establishing the rank--frequency regularity in English. Read for historical grounding; the data analysis is crude by modern standards but the observation is real.

  \item \textbf{S.~T. Piantadosi}, ``Zipf's word frequency law in natural language: a critical review and future directions,'' \emph{Psychon.\ Bull.\ Rev.} (2014) \cite{Piantadosi2014} --- A systematic critical review of explanations for Zipf's law, arguing that matching the rank--frequency plot is weak evidence for any model. The benchmark methodological reference before engaging any proposed mechanism.

  \item \textbf{S.~Goldwater, T.~L. Griffiths, and M.~Johnson}, ``Producing power-law distributions and damping word frequencies with two-stage language models,'' \emph{J.\ Mach.\ Learn.\ Res.} (2011) \cite{GoldwaterGriffithsJohnson2011} --- Derives Zipf-like frequency distributions from a hierarchical Pitman--Yor process and validates the model on English and Chinese corpora. The standard probabilistic explanation in the NLP literature.

  \item \textbf{V.~Berman}, ``Random text, Zipf's law, critical length, and implications for large language models,'' \emph{arXiv:2511.17575} (2025) \cite{Berman2025} --- Derives the Zipf exponent from a combinatorial vocabulary-growth argument, without Bayesian priors. Connects the rank--frequency exponent directly to the observable rate at which new word types appear.

  \item \textbf{N.~Mikhaylovskiy}, ``Zipf's law and the temperature window in transformer language models,'' \emph{Findings of EMNLP} (2025) \cite{Mikhaylovskiy2025b} --- Identifies a finite temperature window within which transformer outputs satisfy Zipf's law, linking the rank--frequency exponent to the phase transitions of Chapter~15.

  \item \textbf{E.~J. Michaud, Z.~Liu, U.~Girit, and M.~Tegmark}, ``The quantization model of neural scaling laws,'' \emph{Proc.\ NeurIPS} (2023) \cite{MichaudEtAl2023} --- Proposes a discrete-quanta model of neural representations in which power-law frequency distributions and neural scaling laws share a common origin. Gives a mechanistic neural-network interpretation of the two-stage Pitman--Yor hierarchy.
\end{itemize}

\section*{Exercises}
\label{ch16:exercises}

\begin{exercise}
\label{ch16:ex1}
Let $p_r = C r^{-s}$ for $r = 1, 2, \ldots, V$ with $C$ chosen so that $\sum_{r=1}^{V} p_r = 1$. For a corpus of $n$ i.i.d. draws from this distribution, compute the expected number of distinct types seen after $n$ draws as a function of $n$, $V$, and $s$. For $s = 1$ and $V \to \infty$, show that this expected count grows as $\Theta(C \log n)$. Compare with Berman's prediction that $V(n) \sim n^{\beta}$ for some $\beta$; for what $\beta$ are the two consistent?
\end{exercise}

\begin{exercise}
\label{ch16:ex2}
Consider the Pitman--Yor Chinese restaurant process with parameters $a = 0.5$, $\theta = 1$. Simulate $n = 10000$ token draws. Plot the empirical rank--frequency distribution on log--log axes and estimate the slope. Compare your empirical estimate with the theoretical prediction $s = 1 + 1/a = 3$. How many tokens are needed before the power-law regime is clearly visible? What does this say about the difficulty of distinguishing power-law tails from, say, a log-normal distribution with comparable data?
\end{exercise}
